\documentclass[11pt]{article}
\usepackage{explorer}
\newcommand{\ExplorerLanguage}{EN}
\newcommand{\ExplorerNumber}{09}
\newcommand{\ExplorerKind}{Lesson and method}
\begin{document}
\ExplorerTitle{Advanced high school · geometry}{Rotations and complex numbers}
\ExplorerSection{Before you begin}
Algebraic form, modulus and coordinates.\par
\ExplorerSection{Ideas and method}
The complex number \(z=x\)+iy represents (x,y). Multiplying by i gives \(iz=-y+ix\), which represents (−y,x): a counterclockwise quarter-turn about the origin.\par
Multiplication by −1 gives a half-turn, and by −i a clockwise quarter-turn. They preserve distances from the origin because |\(i|=1\).\par
For rotation about centre a, translate to the origin, rotate and translate back: \(z'=a+i(z-a)\). This counterclockwise rotation preserves distances between all points.\par
\ExplorerSection{Worked example}
With \(a=1+i\) and \(z=3+2i, z-a=2+i\) and \(i(z-a)=-1+2i\). The image is \(z'=3i\). Relative vectors (2,1) and (−1,2) are perpendicular and both have norm \(\sqrt{5}\).\par
\begin{center}\begin{tikzpicture}[scale=.85]\draw[->,gray](-.5,0)--(3,0) node[right]{$x$};\draw[->,gray](0,-.5)--(0,3) node[above]{$y$};\draw[mint,thick](2,0)--(2,2)--(0,2)--(0,0)--cycle;\fill[violet](1,1)circle(2pt)node[right]{$a=1+i$};\foreach \x/\y/\z in {2/0/2,2/2/2+2i,0/2/2i,0/0/0}{\fill[mint](\x,\y)circle(2pt)node[below left]{$\z$};}\end{tikzpicture}\end{center}
\ExplorerSection{Explore}
In the complex-number lab, move z and compare z with iz. Check equal moduli and perpendicular vectors.\par
\end{document}
